\subsection{Computable certificates under partial observation}
\label{sec:computablepartial}

The class gap in \cref{thm:classgap} is existential. A finite-observation
specialization admits a different, constructive route: compute a lower bound
on the unrestricted partially observed optimum and an upper bound on the
cost of an implementable controller. Their difference is a global certificate.
Belief-state dynamic programming and its piecewise-linear value functions are
classical \citep{smallwood1973}; computable bounds and finite-horizon
point-based planning are developed by \citet{lovejoy1991} and
\citet{walravenspaan2019}. The construction below combines those foundations
with the execution ledger, an observable finite-memory signature policy, and
the model-error certificate of \cref{cor:bookcertificate}.

\begin{assumption}[Finite observations and a binary hidden state]
\label{ass:finiteobservation}
There are $K$ decision stages, a finite observable state space $\mathcal X$,
a hidden state $h\in\{0,1\}$, and finite action and observation alphabets.
The nonempty feasible set $A_k(x)$ depends only on the observable state,
which retains every ledger, reservation and session variable needed for
feasibility. The joint controlled kernel is
\[
 M_{k,x,a,o}(h,h')\ge0,\qquad
 \sum_{o,h'}M_{k,x,a,o}(h,h')=1.
\]
An outcome $o$ determines the next observable state $x'=F_k(x,a,o)$ and
the stage charge $c_k(x,a,o)$. A terminal charge $c_K(x,h)$ is allowed.
Actions may change both the hidden transition and the observation law.
For each stage $j$, all remaining total losses lie in a common interval of
width $S_j$ over states and feasible continuation policies. All primitive
arrays, charges and these bounds are supplied. The initial state $x_0$ and
belief $b_0=\Pp(h_0=1)$ are known.
\end{assumption}

The analytic statements allow real data. Algorithmic claims require rational
primitives and grid nodes, or validated enclosures together with certified
comparisons for the controller's memory update.

The observation alphabet is the information actually available to this
controller. If a marked continuous-time model is reduced to these symbols,
discarding timestamps or reports changes the information problem unless
sufficiency is proved. No equivalence with a richer observation filtration
is assumed. Section~\ref{sec:summarycertificate} supplies a separate
certificate for that richer comparison. Binary hidden liquidity is the explicit computable case below;
it does not exhaust the finite hidden-state setting.

For $b\in[0,1]$, put $\boldsymbol b=(1-b,b)$ and define
\begin{align}\label{eq:finitebeliefkernel}
 p_b^a(o)&=\sum_{h,h'}\boldsymbol b_h M_{k,x,a,o}(h,h'),&
 \Phi_b^a(o)&=\frac{\sum_h\boldsymbol b_h M_{k,x,a,o}(h,1)}{p_b^a(o)},\\
 r_k(x,b,a)&=\sum_o p_b^a(o)c_k(x,a,o).\nonumber
\end{align}
Zero-probability outcomes contribute zero and may be assigned any posterior.
The optimal value $V_k(x,b)$ minimizes expected remaining cost over all
policies using the declared observation history, including their own actions.

\begin{lemma}[Concavity and a computable belief modulus]\label{lem:beliefmodulus}
For fixed $k,x$, $V_k(x,\cdot)$ is concave, piecewise linear and
$S_k$-Lipschitz on $[0,1]$.
\end{lemma}
\begin{proof}
A fixed deterministic continuation tree has conditional costs $v_0,v_1$,
so its expected cost is $(1-b)v_0+bv_1$. The common feasible action masks
give the same finite family of trees for every $b$. Their minimum is concave
and piecewise linear. The range assumption gives $|v_1-v_0|\le S_k$ for
every tree; taking minima preserves that Lipschitz bound. Behavioral
randomization cannot improve the minimum of these expected costs.
\end{proof}

\subsubsection*{Two finite backward recursions}

Let $G_m=\{0,1/m,\ldots,1\}\cup\{b_0\}$, and let $Q_m b$ be a
nearest node, with a fixed tie rule. Write $I_m f$ for linear interpolation
of node values: if $b=(1-\theta)g_i+\theta g_{i+1}$, then
$(I_m f)(b)=(1-\theta)f(g_i)+\theta f(g_{i+1})$.
Adjacent nodes have distance at most $1/m$; inserting the initial prior
therefore preserves the mesh estimates without restricting its value.
The nearest-node planning problem is the finite Markov decision process
\begin{align}\label{eq:nearestbeliefdp}
 W_K(x,g)&=(1-g)c_K(x,0)+g\,c_K(x,1),\nonumber\\
 W_k(x,g)&=\min_{a\in A_k(x)}\left\{
 r_k(x,g,a)+\sum_o p_g^a(o)
 W_{k+1}\bigl(F_k(x,a,o),Q_m\Phi_g^a(o)\bigr)\right\}.
\end{align}
Let $a_k^m(x,g)$ be a minimizing action. Independently compute
\begin{align}\label{eq:belieflowerdp}
 \underline V_K^m(x,g)&=(1-g)c_K(x,0)+g\,c_K(x,1),\nonumber\\
 \underline V_k^m(x,g)&=\min_{a\in A_k(x)}\left\{
 r_k(x,g,a)+\sum_o p_g^a(o)
 (I_m\underline V_{k+1}^m)\bigl(F_k(x,a,o),\Phi_g^a(o)\bigr)\right\}.
\end{align}
Interpolation here gives a lower relaxation for cost minimization. It is
not the deployed controller's state update.

The deployed controller starts at $\gamma_0=b_0$, selects
$a_k^m(x_k,\gamma_k)$, and updates its observable memory by
\begin{equation}\label{eq:beliefmemory}
 \gamma_{k+1}=Q_m\Phi_{\gamma_k}^{a_k}(o_{k+1}).
\end{equation}
This recursively rounded belief need not equal the true posterior conditional
on the entire observed history. The physical hidden state always evolves
according to $M$; it is never redrawn from $\gamma_k$ in deployment.

\subsubsection*{Evaluation in the physical model}

For any feasible finite-memory action probabilities $\sigma_k(a\mid x,\gamma)$
and the update \eqref{eq:beliefmemory}, evaluate on the augmented state
$(x,\gamma,h)$:
\begin{align}\label{eq:truecontrollerevaluation}
 E_K(x,\gamma,h)&=c_K(x,h),\nonumber\\
 E_k(x,\gamma,h)&=\sum_{a\in A_k(x)}\sigma_k(a\mid x,\gamma)
 \sum_{o,h'}M_{k,x,a,o}(h,h')\bigl\{c_k(x,a,o)\nonumber\\
 &\hspace{23mm}+E_{k+1}\bigl(F_k(x,a,o),Q_m\Phi_\gamma^a(o),h'\bigr)\bigr\},\\
 U_m(\sigma)&=(1-b_0)E_0(x_0,b_0,0)+b_0E_0(x_0,b_0,1).\nonumber
\end{align}
The hidden coordinate is integrated out offline. Neither the online action
nor its memory update observes $h$. Evaluation may also use a different
physical kernel while retaining the planner's filter in the memory update.

\begin{theorem}[A computable global execution certificate]
\label{thm:computablepartial}
Under \cref{ass:finiteobservation}, write $V^*=V_0(x_0,b_0)$. For every
feasible controller evaluated by \eqref{eq:truecontrollerevaluation},
\begin{equation}\label{eq:computablegap}
 \underline V_0^m(x_0,b_0)\le V^*\le U_m(\sigma),\qquad
 0\le J_\sigma-V^*\le
 \mathcal G_m(\sigma):=U_m(\sigma)-\underline V_0^m(x_0,b_0).
\end{equation}
All quantities on the right are obtained by finite backward induction.
An outward lower enclosure of $\underline V_0^m$ and outward upper enclosure
of $U_m$ give a valid numerically enclosed certificate. The result holds
for any candidate table; its validity does not require a local optimizer
to have found a global minimum.
\end{theorem}
\begin{proof}
At $K$ the grid terminal values are exact. If
$\underline V_{k+1}^m(x,g_i)\le V_{k+1}(x,g_i)$ at all nodes, concavity
from \cref{lem:beliefmodulus} gives
$I_m\underline V_{k+1}^m(x,b)\le V_{k+1}(x,b)$ for every belief.
Nonnegative observation probabilities and minimization preserve the inequality
in the backward step. This proves the lower bound. Conditioning on the
actual hidden state and the finite controller memory gives
\eqref{eq:truecontrollerevaluation}, hence $U_m(\sigma)=J_\sigma$.
Feasibility places this controller in the class defining $V^*$ and proves
the upper bound. Subtraction gives the certificate. Replacing either endpoint
by an enclosure in the indicated direction can only enlarge it.
\end{proof}

This certificate closes both an optimization gap and a representation gap
relative to the full history-policy optimum of this model. Successive mesh
values need not yield monotonically improving deployed policies. A stopping
rule uses the computed gap of the retained best certified controller, not
the difference between two grid values.

\begin{proposition}[Explicit mesh and numerical optimization bounds]
\label{prop:partialmeshrate}
Set
\[
 \varepsilon_m=\frac1{2m}\sum_{j=1}^K S_j.
\]
For the exact nearest-node minimizer $\sigma^m$,
\begin{equation}\label{eq:partialmeshrate}
 J_{\sigma^m}-V^*\le2\varepsilon_m,\qquad
 \mathcal G_m(\sigma^m)\le3\varepsilon_m.
\end{equation}
If a computed table has global gap at most $\tau_m$ in the finite problem
\eqref{eq:nearestbeliefdp}, add $\tau_m$ to both bounds. One computable
choice is $\tau_m=\sum_k\tilde e_k$, where $\tilde e_k$ bounds, at every grid state,
the selected action's excess over the minimum of its true Bellman action
values. Enclosures
$\widetilde Q_k^-(z,a)\le\widetilde Q_k(z,a)\le\widetilde Q_k^+(z,a)$ give
\begin{equation}\label{eq:partialnumericalgap}
 \tilde e_k=\max_z\left[\widetilde Q_k^+(z,\widehat a_k(z))-
                         \min_{a\in A_k(z)}\widetilde Q_k^-(z,a)\right]_+.
\end{equation}
Here the action values use the optimal continuation of the finite model,
enclosed recursively by backward induction.
\end{proposition}
\begin{proof}
Nearest rounding changes a binary belief by at most $1/(2m)$ in total
variation. For interpolation between neighboring nodes, the mean distance
of its two nodes from $b$ is
$2\theta(1-\theta)(g_{i+1}-g_i)\le1/(2m)$.
Using \cref{lem:beliefmodulus} in the two backward recursions gives,
at grid nodes,
\[
 |W_k-V_k|\le\frac1{2m}\sum_{j=k+1}^K S_j,\qquad
 0\le V_k-\underline V_k^m\le\frac1{2m}\sum_{j=k+1}^K S_j.
\]
The terminal interpolation is in fact exact; retaining $S_K$ is a valid
conservative bound.

To compare the deployed controller with the nearest-node planning model,
introduce an artificial hidden-state resampling after each observation.
Given the preceding artificial history, the pre-resampling posterior is
$\Phi_\gamma^a(o)$; replace it by $Q_m\Phi_\gamma^a(o)$. This defines the
finite Markov model \eqref{eq:nearestbeliefdp}. Starting from the physical
model, introduce resamplings one at a time in forward order: the $j$th
hybrid has the first $j$ resamplings and no later ones. Before the next
replacement, the common prefix therefore has all preceding resamplings,
so the stated posterior is its correct conditional law. Any fixed
continuation rule, with no later resampling in either compared hybrid, has conditional
remaining-cost span at most $S_j$. Changing this distribution changes its
expected cost by at most $S_j/(2m)$. Summing gives a uniform objective
difference at most $\varepsilon_m$ for every common history policy.
This argument compares unconditional policy costs; it needs no uniform
Lipschitz bound on Bayes' formula after a rare observation.

The artificial model has the observable sufficient state $(x,\gamma)$,
so its optimum over all history policies is $W_0(x_0,b_0)$. Comparing the
true and artificial optima and their selected policy gives
$J_{\sigma^m}-V^*\le2\varepsilon_m$. Adding the lower-relaxation defect
proves the factor three. A finite-model optimization gap contributes
$\tau_m$ once. Finally, the finite-horizon performance-difference recursion
bounds that gap by $\sum_k\tilde e_k$; \eqref{eq:partialnumericalgap} encloses each
selected-action excess.
\end{proof}

In the resampling comparison, the continuation uses the same observable
memory update and action map in both models. A sampled hidden state is not
revealed. The artificial model exists only for the comparison proof;
\eqref{eq:truecontrollerevaluation} evaluates the physical model.

\begin{corollary}[Finite-memory signature implementation and model transfer]
\label{cor:partialsignature}
Append an observable summary channel which starts at zero and at decision
$k$ has value $e_{(x_k,\gamma_k)}$. Linear interpolation between its update
nodes gives a finite-variation path whose first signature level equals this
current one-hot vector. Give $a_k^m(x,\gamma)$ score $\log w$, $w\ge1$,
and every other admissible action score zero, retaining the exact mask.
This is a depth-one signature-softmax controller $\sigma^{m,w}$ on the
specified augmented observation path. With $m_A=|A|$,
\begin{equation}\label{eq:partialsignaturegap}
 \mathcal G_m(\sigma^{m,w})
 \le3\varepsilon_m+\tau_m+
       KS_0\frac{m_A-1}{w+m_A-1}.
\end{equation}
Its usually sharper certificate is the directly computed
$U_m(\sigma^{m,w})-\underline V_0^m$.
If these quantities are computed in an approximate physical model and
$|J_\pi-\widehat J_\pi|\le B_T$ uniformly over the common feasible policies,
then the deployed policy satisfies
\begin{equation}\label{eq:partialtransfer}
 J_{\sigma^{m,w}}-\inf_{\pi\in\Pi}J_\pi
 \le 2B_T+\widehat U_m(\sigma^{m,w})-
                  \widehat{\underline V}_0^m(x_0,b_0).
\end{equation}
For the hypotheses of \cref{cor:bookcertificate}, one may take
$B_T=D_{\rm cost}+S_0(1-e^{-dT})$ when $S_0$ also bounds the approximate
loss range.
\end{corollary}
\begin{proof}
The first signature level is the channel's endpoint minus its zero origin.
Stage-specific linear coefficients therefore reproduce the table scores.
With $m_k(z):=|A_k(z)|\le m_A$ feasible actions, where $z=(x,\gamma)$,
the preferred-action probability is $w/(w+m_k(z)-1)$.
The bound \eqref{eq:partialsignaturegap} uses the monotone worst case
$m_k(z)=m_A$. \Cref{lem:policyperturb} bounds its cost difference from the
deterministic table by the last term in that bound.
Apply \cref{prop:partialmeshrate}. For transfer, insert the approximate
objective on both sides of a policy comparison and use
\cref{thm:computablepartial}; the two model discrepancies contribute $2B_T$.
The filter and table remain fixed when this policy is evaluated in the true
environment.
\end{proof}

The summary is computed from observations, actions and the supplied filter,
not from the hidden regime. Its dimension is
$|\mathcal X||G_m|\le|\mathcal X|(m+2)$, and
the policy has $K|\mathcal X||G_m|$ preferred-action entries. With at most
$B$ outcome branches, the recursions take
$O(K|\mathcal X||G_m||A|B)$ arithmetic operations for a binary hidden state.
The explicit accuracy parameter is the belief mesh and table size. This is
not a dimension-free depth bound for the unaugmented transcript signature,
nor does it make endogenous coefficient optimization convex. A higher-
dimensional hidden-state simplex requires a mesh whose size can grow
combinatorially. Unknown primitive-model error still requires evidence;
the computed planning gap alone cannot certify an unvalidated market model.
